\section{AIME II 2007}
        \begin{enumerate}
        	\item (\textit{AIME II 2007}) A mathematical organization is producing a set of commemorative license plates. Each plate contains a sequence of five characters chosen from the four letters in AIME and the four digits in $2007$. No character may appear in a sequence more times than it appears among the four letters in AIME or the four digits in $2007$. A set of plates in which each possible sequence appears exactly once contains N license plates. Find N/10.


 
 \textbf{Answer: }372



\textbf{Solution 1}

There are 7 different characters that can be picked, with 0 being the only number that can be repeated twice.

 
                \begin{itemize}
                    \item If $0$ appears 0 or 1 times amongst the sequence, there are $\frac{7!}{(7-5)!} = 2520$ sequences possible.\item If $0$ appears twice in the sequence, there are ${5\choose2} = 10$ places to place the $0$s. There are $\frac{6!}{(6-3)!} = 120$ ways to place the remaining three characters. In total, that gives us $10 \cdot 120 = 1200$.
                \end{itemize}


                Thus, $N = 2520 + 1200 = 3720$, and $\frac{N}{10} = \boxed{372}$.

 Note: We could also count cases by how many 0's there are, like No 0's, 1 0's, and 2 0's. We get the same answer.

 





\textbf{Solution 2}

Another approach is to first falsely assume that the zeros are distinguishable and then correct the over-counting.

 We can name the zeros $0_1$ and $0_2$. The first character has 8 choices - $A,I,M,E,2,0_1,0_2,$ and $7$. The second character has $7$ choices, the third has $6$ choices, the fourth has $5$ choices, and the fifth has $4$ choices. Multiplying them all together, there are $6,720$ arrangements where the zeros are distinguishable.

 However, some cases with at most 2 zeros are over-counted. We can tackle this by going case by case for the number of zeros. When there are no zeros there will be no arrangements over-counted. 

 When there is exactly $1$ zero, all the cases are over-counted. This is because there will be another case with the exact same characters but the zero will be switched with the other zero. The number of ways to place the zero is $5$, there are $2$ zeros to choose from, and the remaining spots have $6 \cdot 5 \cdot 4 \cdot 3$ ways to place the non-zero characters. Thus, there are $5 \cdot 2 \cdot 6 \cdot 5 \cdot 4 \cdot 3 = 3,600$ arrangements where there is exactly one zero. However, we need to remove half of these or $1,800$ from the total yielding in $4,920$ arrangements left.

 Cases with two zeros would also have to be corrected as the zeros are interchangeable and you can switch the ordering of them. The number of ways to place $0_1$ is $5$, the number of ways to place $0_2$ is $4$, and the remaining places have $6 \cdot 5 \cdot 4$ ways to place the remaining characters. In total, there are $5 \cdot 4 \cdot 6 \cdot 5 \cdot 4 = 2,400$ arrangements with exactly two zeros. However, we need to remove half of these or $1,200$ from the total yielding in $3,720$ arrangements left.

 The max number of zeros is two so we have gone through all the cases and removed the over-counted arrangements thus giving, $N = 3,720$ and $\frac{N}{10} = \boxed{372}$.

 





\textbf{Solution 3}

Another way to solve this is through brute force casework.

 We can have 1, 2, 3, or 4 letters in the sequence, each from AIME. 

 \textbf{Case 1: 4 letters}

 Here, we have a string in the base form $AIMEx$ where $x \in \{0,2,7 \}$. That is three ways times $5!$ permutations, being 360 cases.

 \textbf{Case 2: 3 letters}

 Assume the string has base form with no zeroes. Then, it takes the form $AIM27$, which has $5!$ permutations, times the 4 ways to choose the three letters, $5! \cdot 4$. Say it has one zero. Then, it takes the form $AIM0x$ for $x \in \{2,7 \}$, which is $2 \cdot 4 \cdot  5!$. Otherwise, the string has two zeroes, and that is $AIM00$, which is $4 \cdot \frac{5!}{2!}$. Summing, we have 1680 cases.

 \textbf{Case 3: 2 letters}

 No base form with no zeroes exists. Therefore, the base form has one zero. Then it takes the form $AI207$, which is $6 \cdot 5!$ ways. Otherwise, it has two zeroes, and the base form is $AI00x$ with $x \in \{2,7 \}$. Then, there are $\frac{5!}{2!} \cdot 2 \cdot 6$ cases here. Summing, we have 1440 cases.

 \textbf{Case 4: 1 letter}

 The number 2007 must be in the string. Otherwise, we have 4 choices for the letter, and $\frac{5!}{2!}$ total arrangements. Then, we have $\frac{5!}{2!} \cdot 4 = 240$ cases.

 Every subcase here ends in a 0, so when the total count $N$ is divided by 10, we just remove the 0's in the subcases. The requested $N$ is then $24 + 144 + 168+36 = \boxed{372}$.

 




	\item (\textit{AIME II 2007}) Find the number of ordered triples $(a,b,c)$ where $a$, $b$, and $c$ are positive \href{https://artofproblemsolving.com/wiki/index.php?title=Integer}{integers}, $a$ is a \href{https://artofproblemsolving.com/wiki/index.php?title=Factor}{factor} of $b$, $a$ is a factor of $c$, and $a+b+c=100$.


 
 \textbf{Answer: }200



\textbf{Solution 1}

Denote $x = \frac{b}{a}$ and $y = \frac{c}{a}$. The last condition reduces to $a(1 + x + y) = 100$. Therefore, $1 + x + y$ is equal to one of the 9 factors of $100 = 2^25^2$.

 Subtracting the one, we see that $x + y = \{0,1,3,4,9,19,24,49,99\}$. There are exactly $n - 1$ ways to find pairs of $(x,y)$ if $x + y = n$. Thus, there are $0 + 0 + 2 + 3 + 8 + 18 + 23 + 48 + 98 = \boxed{200}$ solutions of $(a,b,c)$.

 Alternatively, note that the sum of the divisors of $100$ is $(1 + 2 + 2^2)(1 + 5 + 5^2)$ (notice that after distributing, every divisor is accounted for). This evaluates to $7 \cdot 31 = 217$. Subtract $9 \cdot 2$ for reasons noted above to get $199$. Finally, this changes $1 \Rightarrow -1$, so we have to add one to account for that. We get $\boxed{200}$.

 





\textbf{Solution 2 (similar to 1, different ending)}

Let $b = ak$ and $c = aj$ for $k,j \in \mathbb{Z}^+ \neq 0$, as $b,c \in \mathbb{Z}^+$. Then, factoring out $a$, we have $a(1+x+y) = 100$. $a$ is an integer, and so is $1+x+y$ by definition, so that implies $a \ | \ 100$. The possible choices for $a$ are $1, 2, 4, 5, 10, 20, 25, 50, 100$. Each of those gives options $k+j = 0,1,3,4,9,19,24,49,99$ respectively. Negate 0 and 1 as $k$ and $j$ are positive integers not equal to 0. While we can attempt to find through exclusion stars and bars the number of ways as $\binom{x-1}{y-1}$, coming to the realization that $y-1$ will always be 1 (as we have 2 variables $k,j$), and continuing by solution 1, we could also constructively overcount. By stars and bars, each two equations will have integer solutions, inclusive 0, as $\binom{x+y-1}{y-1}$. Note that $y=2$ forces the bottom of the binomial as 1, making it expendable and therefore we just consider $x+y-1$. Here, $y-1 = 1$ so we just consider $k+1$ cases total, so we sum over ALL terms and obtain $100+50+25+20+10+5+4 = 214$. However, 14 of these cases are invalid, particularly, when $k+j = 3,4,9,19,24,49,99$, and one of $k$ or $j$ is 0. There are two choice for which variable is 0, and 7 total $k+j$ values, which is $2 \cdot 7 = 14$ of the total that do not work, so the answer is $214 - 14 = \boxed{200}$.

 




	\item (\textit{AIME II 2007}) \href{https://artofproblemsolving.com/wiki/index.php?title=Square}{Square} $ABCD$ has side length $13$, and \href{https://artofproblemsolving.com/wiki/index.php?title=Point}{points} $E$ and $F$ are exterior to the square such that $BE=DF=5$ and $AE=CF=12$. Find $EF^{2}$.


 
                    \begin{figure}[H]
                        \centering
                        \includegraphics[width=0.5\textwidth]{Problem/AIME_II_2007_3_Problem_Figure_1.png}
                    \end{figure}
                    
 \textbf{Answer: }578



\textbf{Solution 1}

Let $\angle FCD = \alpha$, so that $FB = \sqrt{12^2 + 13^2 + 2\cdot12\cdot13\sin(\alpha)} = \sqrt{433}$. By the diagonal, $DB = 13\sqrt{2}, DB^2 = 338$. 

 \textbf{The sum of the squares of the sides of a parallelogram is the sum of the squares of the diagonals.}
 \[EF^2 = 2\cdot(5^2 + 433) - 338 = 578.\]

 





\textbf{Solution 2}

Extend $\overline{AE}, \overline{DF}$ and $\overline{BE}, \overline{CF}$ to their points of intersection. Since $\triangle ABE \cong \triangle CDF$ and are both $5-12-13$ \href{https://artofproblemsolving.com/wiki/index.php?title=Right\_triangle}{right triangles}, we can come to the conclusion that the two new triangles are also congruent to these two (use \href{https://artofproblemsolving.com/wiki/index.php/Congruent\_(geometry)\#ASA\_Congruence}{ASA}, as we know all the sides are $13$ and the angles are mostly complementary). Thus, we create a \href{https://artofproblemsolving.com/wiki/index.php?title=Square}{square} with sides $5 + 12 = 17$.

 
                    \begin{figure}[H]
                        \centering
                        \includegraphics[width=0.5\textwidth]{Solution/AIME_II_2007_3_Solution_Figure_1.png}
                    \end{figure}
                    

 $\overline{EF}$ is the diagonal of the square, with length $17\sqrt{2}$; the answer is $EF^2 = (17\sqrt{2})^2 = 578$.

 





\textbf{Solution 3}

A slightly more analytic/brute-force approach:

 
                    \begin{figure}[H]
                        \centering
                        \includegraphics[width=0.5\textwidth]{Solution/AIME_II_2007_3_Solution_Figure_2.PNG}
                    \end{figure}
                    

 Drop perpendiculars from $E$ and $F$ to $I$ and $J$, respectively; construct right triangle $EKF$ with right angle at K and $EK || BC$. Since $2[CDF]=DF*CF=CD*JF$, we have $JF=5\times12/13 = \frac{60}{13}$. Similarly, $EI=\frac{60}{13}$. Since $\triangle DJF \sim \triangle DFC$, we have $DJ=\frac{5JF}{12}=\frac{25}{13}$.

 Now, we see that $FK=DC-(DJ+IB)=DC-2DJ=13-\frac{50}{13}=\frac{119}{13}$. Also, $EK=BC+(JF+IE)=BC+2JF=13+\frac{120}{13}=\frac{289}{13}$. By the Pythagorean Theorem, we have $EF=\sqrt{\left(\frac{289}{13}\right)^2+\left(\frac{119}{13} \right)^2}=\frac{\sqrt{(17^2)(17^2+7^2)}}{13}$$=\frac{17\sqrt{338}}{13}=\frac{17(13\sqrt{2})}{13}=17\sqrt{2}$. Therefore, $EF^2=(17\sqrt{2})^2=578$.

 





\textbf{Solution 4}

Based on the symmetry, we know that $F$ is a reflection of $E$ across the center of the square, which we will denote as $O$. Since $\angle BEA$ and $\angle AOB$ are right, we can conclude that figure $AOBE$ is a cyclic quadrilateral. Pythagorean Theorem yields that $BO=AO=\frac{13\sqrt{2}}{2}$. Now, using Ptolemy's Theorem, we get that 
 \[AO\cdot BE + BO\cdot AE = AB\cdot AO\]
 \[\frac{13\sqrt{2}}{2}\cdot 5+\frac{13\sqrt{2}}{2}\cdot 12 = 13\cdot OE\]
 \[OE=\frac{17\sqrt{2}}{2}\]Now, since we stated in the first step that $F$ is a reflection of $E$ across $O$, we can say that $EF=2EO=17\sqrt{2}$. This gives that 
 \[EF^2=(17\sqrt{2})^2=578\] AWD with this bash solution

 





\textbf{Solution 5 (Ptolemy's Theorem)}

Drawing $EF$, it clearly passes through the center of $ABCD$. Letting this point be $P$, we note that $AEBP$ and $CFDP$ are congruent cyclic quadrilaterals, and that $AP=BP=CP=DP=\frac{13}{\sqrt{2}}.$ Now, from Ptolemy's, $13\cdot EP=\frac{13}{\sqrt{2}}(12+5)\implies EP=\frac{17\sqrt{2}}{2}$. Since $EF=EP+FP=2\cdot EP$, the answer is $(17\sqrt{2})^2=\boxed{578}.$

 





\textbf{Solution 6 (Coordinate Bash)}

Place the vertices of the square as follows:\begin{align}D &= (0,\;0),\\C &= (13,\;0),\\A &= (0,\;13),\\B &= (13,\;13).\end{align}

 Compute the coordinates of points E and F using the distance formula.Place $E$ at $(x,y)$ and $F$ at $(a,b)$.

 \begin{aligned}x\textasciicircum{}2+(y-13)\textasciicircum{}2 = 144,\\(x-13)\textasciicircum{}2+(y-13)\textasciicircum{}2 = 25\end{aligned}

 \begin{aligned}a\textasciicircum{}2+b\textasciicircum{}2=25,\\(a-13)\textasciicircum{}2+b\textasciicircum{}2=144.\end{aligned}

 Solving and taking appropriate solutions to get:$E = (144/13,229/13)$, and $F = (25/13, -60/13)$

 Computing distance between $E$ and $F$ = $17\sqrt{2}$

 
 \[EF^2 = 578\]

 





\textbf{Solution 7 (Trig Bash)}

We first see that the whole figure is symmetrical and reflections across the center that we will denote as $O$ bring each half of the figure to the other half. Thus we consider a single part of the figure, namely $EO.$

 First note that $\angle BAO = 45^{\circ}$ since $O$ is the center of square $ABCD.$ Also note that $\angle EAB = \arccos{\left(\frac{12}{13}\right)}$ or $\arcsin{\left(\frac{5}{13}\right)}.$ Finally, we know that $AO =\frac{13\sqrt{2}}{2}.$ Now we apply laws of cosines on $\bigtriangleup AEO.$ 

 We have $EO^2 = 12^2 + (\frac{13\sqrt{2}}{2})^2 - 2 \cdot 12 \cdot \left(\frac{13\sqrt{2}}{2}\right) \cdot \cos{\angle EAO}.$ We know that $\angle EAO = 45^{\circ} + \arccos{\left(\frac{12}{13}\right)}.$ Thus we have $\cos{\angle EAO} = \cos\left(45^{\circ} + \arccos{\left(\frac{12}{13}\right)}\right)$ which applying the cosine sum identity yields $\cos{45^{\circ}}\cos{\arccos\frac{12}{13}} - \sin{45^{\circ}}\sin\arcsin{\frac{5}{12}} =\frac{12\sqrt{2}}{26} -\frac{5\sqrt{2}}{26} =\frac{7\sqrt{2}}{26}.$ 

 Note that we are looking for $4EO^2$ so we multiply $EO^2 = 12^2 + \left(\frac{13\sqrt{2}}{2}\right)^2 - 2 \cdot 12 \cdot \left(\frac{13\sqrt{2}}{2}\right) \cdot \cos{\angle EAO}$ by $4$ obtaining $4EO^2 = 576 + 338 - 8 \cdot\left(\frac{13\sqrt{2}}{2}\right) \cdot 12 \cdot\frac{7\sqrt{2}}{26} = 576 + 338 - 4  \cdot 12 \cdot  7 = \boxed{578}.$

 





\textbf{Solution 8 (Rectangular Construction)}


                    \begin{figure}[H]
                        \centering
                        \includegraphics[width=0.5\textwidth]{Solution/AIME_II_2007_3_Solution_Figure_3.png}
                    \end{figure}
                    

 Construct a rectangle $A'B'C'D'$ as shown above starting with the top left vertex, moving clockwise about the blue box. Let $EB' = a$, $B'B = c$, and $EA' = b$. Through \textbf{CPCTC} it follows that most other side lengths due to congruence are equal to one of $a,b,c$. Three equations arise, particularly $a+b = 13, a^2 + c^2 = 25, b^2 + c^2 = 144$, which we can easily solve to see $a = \frac{25}{13}$, $b = \frac{144}{13}$, and $c = \frac{60}{13}$ (negate negative solution). Drop the altitude from $E$ down to $C'D'$, call it $G$. It follows that $FGE$ is right with legs $b-a=\frac{144-25}{13}=\frac{119}{13}$ and $13+2c=\frac{289}{13}$, so solving via substitution and Pythagorean yields $EF = \frac{\sqrt{97682}}{13} = \frac{13\sqrt{578}}{13} = \sqrt{578}$. So, $EF^2 = \boxed{578}$.

 




	\item (\textit{AIME II 2007}) The workers in a factory produce widgets and whoosits. For each product, production time is \href{https://artofproblemsolving.com/wiki/index.php?title=Constant}{constant} and identical for all workers, but not necessarily equal for the two products. In one hour, $100$ workers can produce $300$ widgets and $200$ whoosits. In two hours, $60$ workers can produce $240$ widgets and $300$ whoosits. In three hours, $50$ workers can produce $150$ widgets and $m$ whoosits. Find $m$.


 
 \textbf{Answer: }450



\textbf{Solutions}

Suppose that it takes $x$ hours for one worker to create one widget, and $y$ hours for one worker to create one whoosit.

 Therefore, we can write that (note that two hours is similar to having twice the number of workers, and so on):

 $100 = 300x + 200y$

 $2(60) = 240x + 300y$

 $3(50) = 150x + my$

 Solve the system of equations with the first two equations to find that $(x,y) = \left(\frac{1}{7}, \frac{2}{7}\right)$. Substitute this into the third equation to find that $1050 = 150 + 2m$, so $m = \boxed{450}$.

 




	\item (\textit{AIME II 2007}) The \href{https://artofproblemsolving.com/wiki/index.php?title=Graph}{graph} of the \href{https://artofproblemsolving.com/wiki/index.php?title=Equation}{equation} $9x+223y=2007$ is drawn on graph paper with each \href{https://artofproblemsolving.com/wiki/index.php?title=Square}{square} representing one \href{https://artofproblemsolving.com/wiki/index.php?title=Unit\_square}{unit} in each direction. How many of the $1$ by $1$ graph paper squares have interiors lying entirely below the graph and entirely in the first \href{https://artofproblemsolving.com/wiki/index.php?title=Quadrant}{quadrant}?


 
 \textbf{Answer: }888



\textbf{Solution 1}

There are $223 \cdot 9 = 2007$ squares in total formed by the rectangle with edges on the x and y axes and with vertices on the intercepts of the equation, since the \href{https://artofproblemsolving.com/wiki/index.php?title=Intercept&action=edit&redlink=1}{intercepts} of the lines are $(223,0),\ (0,9)$.

 Count the number of squares that the diagonal of the rectangle passes through. Since the two diagonals of a rectangle are \href{https://artofproblemsolving.com/wiki/index.php?title=Congruent}{congruent}, we can consider instead the diagonal $y = \frac{223}{9}x$. This passes through 8 horizontal lines ($y = 1 \ldots 8$) and 222 vertical lines ($x = 1 \ldots 222$). Every time we cross a line, we enter a new square. Since 9 and 223 are \href{https://artofproblemsolving.com/wiki/index.php?title=Relatively\_prime}{relatively prime}, we don’t have to worry about crossing an intersection of a horizontal and vertical line at one time. We must also account for the first square. This means that it passes through $222 + 8 + 1 = 231$ squares.

 The number of non-diagonal squares is $2007 - 231 = 1776$. Divide this in 2 to get the number of squares in one of the triangles, with the answer being $\frac{1776}2 = 888$.

 





\textbf{Solution 2}

Count the number of each squares in each row of the triangle. The \href{https://artofproblemsolving.com/wiki/index.php?title=Intercept&action=edit&redlink=1}{intercepts} of the \href{https://artofproblemsolving.com/wiki/index.php?title=Line}{line} are $(223,0),\ (0,9)$.

 In the top row, there clearly are no squares that can be formed. In the second row, we see that the line $y = 8$ gives a $x$ value of $\frac{2007 - 8(223)}{9} = 24 \frac 79$, which means that $\lfloor 24 \frac 79\rfloor = 24$ unit squares can fit in that row. In general, there are

 $\sum_{i=0}^{8} \lfloor \frac{223i}{9} \rfloor$triangles. Since $\lfloor \frac{223}{9} \rfloor = 24$, we see that there are more than $24(0 + 1 + \ldots + 8) = 24(\frac{8 \times 9}{2}) = 864$ triangles. Now, count the fractional parts. $\lfloor \frac{0}{9} \rfloor = 0, \lfloor \frac{7}{9} \rfloor = 0, \lfloor \frac{14}{9} \rfloor = 1,$ $\lfloor \frac{21}{9} \rfloor = 2, \lfloor \frac{28}{9} \rfloor = 3, \lfloor \frac{35}{9} \rfloor = 3,$ $\lfloor \frac{42}{9} \rfloor = 4, \lfloor \frac{49}{9} \rfloor = 5, \lfloor \frac{56}{9} \rfloor = 6$. Adding them up, we get $864 + 1 + 2 + 3 + 3 + 4 + 5 + 6 = 888$.

 





\textbf{Solution 3}

From \href{https://artofproblemsolving.com/wiki/index.php?title=Pick\%27s\_Theorem}{Pick's Theorem}, $\frac{2007}{2}=\frac{233}{2}-\frac{2}{2}+\frac{2I}{2}$. In other words, $2I=1776$ and I is $888$.

 Do you see why we simply set $I$ as the answer as well? That's because every interior point, if moved down and left one (southwest direction), can have that point and the previous point create a unit square. For example, $(1, 1)$ moves to $(0, 0)$, so the square of points ${(0, 0), (1, 0), (1, 1), (0, 1)}$ is one example. This applies, of course, for $888$ points.

 





\textbf{Solution 4}

We know that the number of squares intersected in an $m\times{n}$ rectangle is $m + n -\mbox{gcd}(m,n)$. So if we apply that here, we get that the number of intersected squares is:

 $9 + 223 - 1 = 231$. 

 Now just subtract that from the total number of squares and divide by 2, since we want the number of squares below the line. 

 So,

 $\frac{2007 - 231}{2} = \frac{1776}{2} = \fbox{888}$

 




	\item (\textit{AIME II 2007}) An integer is called \textit{parity-monotonic} if its decimal representation $a_{1}a_{2}a_{3}\cdots a_{k}$ satisfies $a_{i}<a_{i+1}$ if $a_{i}$ is \href{https://artofproblemsolving.com/wiki/index.php?title=Odd}{odd}, and $a_{i}>a_{i+1}$ if $a_{i}$ is \href{https://artofproblemsolving.com/wiki/index.php?title=Even}{even}. How many four-digit parity-monotonic integers are there? 


 
 \textbf{Answer: }640



\textbf{Solution 1}

Let's set up a table of values. Notice that 0 and 9 both cannot appear as any of $a_1,\ a_2,\ a_3$ because of the given conditions. (Also note that 0 cannot appear as 0 cannot be the first digit of an integer.) A clear pattern emerges.

 For example, for $3$ in the second column, we note that $3$ is less than $4,6,8$, but greater than $1$, so there are four possible places to align $3$ as the second digit.

 <table align="center" style="border:1px solid black;"><tbody><tr><td>Digit  </td><td>1st  </td><td>2nd  </td><td>3rd  </td><td>4th</td></tr><tr><td>0</td><td>0</td><td>0</td><td>0</td><td>64</td></tr><tr><td>1</td><td>1</td><td>4</td><td>16</td><td>64</td></tr><tr><td>2</td><td>1</td><td>4</td><td>16</td><td>64</td></tr><tr><td>3</td><td>1</td><td>4</td><td>16</td><td>64</td></tr><tr><td>4</td><td>1</td><td>4</td><td>16</td><td>64</td></tr><tr><td>5</td><td>1</td><td>4</td><td>16</td><td>64</td></tr><tr><td>6</td><td>1</td><td>4</td><td>16</td><td>64</td></tr><tr><td>7</td><td>1</td><td>4</td><td>16</td><td>64</td></tr><tr><td>8</td><td>1</td><td>4</td><td>16</td><td>64</td></tr><tr><td>9</td><td>0</td><td>0</td><td>0</td><td>64</td></tr></tbody></table>For any number from 1-8, there are exactly 4 numbers from 1-8 that are odd and less than the number or that are even and greater than the number (the same will happen for 0 and 9 in the last column). Thus, the answer is $4^{k-1} \cdot 10 = 4^3\cdot10 = 640$.-yeetdayeet

 





\textbf{Solution 2 (Recursion)}

This problem can be solved via recursion since we are "building a string" of numbers with some condition. We want to create a new string by adding a new digit at the front so we avoid complications ($0$ can't be at the front and no digit is greater than $9$). There are $4$ options to add no matter what$^{\ast}$ $^{\ast}$ $^{\ast}$ so the recursion is $S_n=4S_{n-1}$ where $S_n$ stands for the number of such numbers with $n$ digits. Since $S_1=10$ the answer is $\boxed{640}$.

 $^{\ast}$ $^{\ast}$ $^{\ast}$ divide 0-9 into groups $(E O)(E O)(E O)(E O)(E O)$. If $a_i$ is even, take an odd from a group below $a_i$'s group or an even from a group above its group as $a_{i-1}$. Do the opposite for an odd $a_i$. So for all $a_i$ there is one option for $a_{i-1}$ from each of the $4$ groups that are not $a_i$'s group.

 




	\item (\textit{AIME II 2007}) Given a \href{https://artofproblemsolving.com/wiki/index.php?title=Real\_number}{real number} $x,$ let $\lfloor x \rfloor$ denote the \href{https://artofproblemsolving.com/wiki/index.php?title=Floor\_function}{greatest integer} less than or equal to $x.$ For a certain \href{https://artofproblemsolving.com/wiki/index.php?title=Integer}{integer} $k,$ there are exactly $70$ positive integers $n_{1}, n_{2}, \ldots, n_{70}$ such that $k=\lfloor\sqrt[3]{n_{1}}\rfloor = \lfloor\sqrt[3]{n_{2}}\rfloor = \cdots = \lfloor\sqrt[3]{n_{70}}\rfloor$ and $k$ divides $n_{i}$ for all $i$ such that $1 \leq i \leq 70.$


 Find the maximum value of $\frac{n_{i}}{k}$ for $1\leq i \leq 70.$


 
 \textbf{Answer: }553



\textbf{Solution 1}

For $x = 1$, we see that $\sqrt[3]{1} \ldots \sqrt[3]{7}$ all work, giving 7 integers. For $x=2$, we see that in $\sqrt[3]{8} \ldots \sqrt[3]{26}$, all of the \href{https://artofproblemsolving.com/wiki/index.php?title=Even}{even} numbers work, giving 10 integers. For $x = 3$, we get 13, and so on. We can predict that at $x = 22$ we get 70.

 To prove this, note that all of the numbers from $\sqrt[3]{x^3} \ldots \sqrt[3]{(x+1)^3 - 1}$ divisible by $x$ work. Thus, $\frac{(x+1)^3 - 1 - x^3}{x} + 1  = \frac{3x^2 + 3x + 1 - 1}{x} + 1 = 3x + 4$ (the one to be inclusive) integers will fit the conditions. $3k + 4 = 70 \Longrightarrow k = 22$.

 The maximum value of $n_i = (x + 1)^3 - 1$. Therefore, the solution is $\frac{23^3 - 1}{22} = \frac{(23 - 1)(23^2 + 23 + 1)}{22} = 529 + 23 + 1 = 553$.

 





\textbf{Solution 2}

Obviously $k$ is positive. Then, we can let $n_1$ equal $k^3$ and similarly let $n_i$ equal $k^3 + (i - 1)k$. 

 The wording of this problem (which uses "exactly") tells us that $k^3+69k<(k+1)^3 = k^3 + 3k^2 + 3k+1 \leq k^3+70k$. Taking away $k^3$ from our inequality results in $69k<3k^2+3k+1\leq 70k$. Since $69k$, $3k^2+3k+1$, and $70k$ are all integers, this inequality is equivalent to $69k\leq 3k^2+3k<70k$. Since $k$ is positive, we can divide the inequality by $k$ to get $69 \leq 3k+3 < 70$. Clearly the only $k$ that satisfies is $k=22$. 

 Then, $\frac{n_{70}}{k}=k^2+69=484+69=\boxed{553}$ is the maximum value of $\frac{n_i}{k}$. (Remember we set $n_i$ equal to $k^3 + (i - 1)k$!)

 


 \href{https://mathproblemsolvingskills.wordpress.com/}{MathProblemSolvingSkills.com}

 


 




	\item (\textit{AIME II 2007}) A \href{https://artofproblemsolving.com/wiki/index.php?title=Rectangle}{rectangular} piece of paper measures 4 units by 5 units. Several \href{https://artofproblemsolving.com/wiki/index.php?title=Line}{lines} are drawn \href{https://artofproblemsolving.com/wiki/index.php?title=Parallel}{parallel} to the edges of the paper. A rectangle determined by the \href{https://artofproblemsolving.com/wiki/index.php?title=Intersection}{intersections} of some of these lines is called \textit{basic} if 


 
            \begin{enumerate}[label=(\roman*)]
                \item  all four sides of the rectangle are segments of drawn line segments, and
\item  no \href{https://artofproblemsolving.com/wiki/index.php?title=Segment}{segments} of drawn lines lie inside the rectangle.
            \end{enumerate}
            Given that the total length of all lines drawn is exactly 2007 units, let $N$ be the maximum possible number of basic rectangles determined. Find the \href{https://artofproblemsolving.com/wiki/index.php?title=Remainder}{remainder} when $N$ is divided by 1000.


 
 \textbf{Answer: }896



\textbf{Solution 1}

Denote the number of horizontal lines drawn as $x$, and the number of vertical lines drawn as $y$. The number of basic rectangles is $(x - 1)(y - 1)$. $5x + 4y = 2007 \Longrightarrow y = \frac{2007 - 5x}{4}$. Substituting, we find that $(x - 1)\left(-\frac 54x + \frac{2003}4\right)$.

 \href{https://artofproblemsolving.com/wiki/index.php?title=FOIL}{FOIL} this to get a quadratic, $-\frac 54x^2 + 502x - \frac{2003}4$. Use $\frac{-b}{2a}$ to find the maximum possible value of the quadratic: $x = \frac{-502}{-2 \cdot \frac 54} = \frac{1004}5 \approx 201$. However, this gives a non-integral answer for $y$. The closest two values that work are $(199,253)$ and $(203,248)$. 

 We see that $252 \cdot 198 = 49896 > 202 \cdot 247 = 49894$. The solution is $\boxed{896}$.

 





\textbf{Solution 2}

We realize that drawing $x$ vertical lines and $y$ horizontal lines, the number of basic rectangles we have is $(x-1)(y-1)$. The easiest possible case to see is $223$ vertical and $223$ horizontal lines, as $(4+5)223 = 2007$. Now, for every 4 vertical lines you take away, you can add 5 horizontal lines, so you basically have the equation $(222-4x)(222+5x)$ maximize. 

 Expanded, this gives $-20x^{2}+222x+222^{2}$. From $-\frac{b}{2a}$ you get that the vertex is at $x=\frac{111}{20}$. This is not an integer though, so you see that when $x=5$, you have $-20*25+222*5+222^{2}$ and that when x=6, you have $-20*36+222*6+222^{2}$. $222 > 20*11$, so the maximum integral value for x occurs when $x=6$. Now you just evaluate $-20*36+222*6+222^{2}\mod 1000$ which is ${896}$.

 





\textbf{Solution 3 (similar to solution 2)}

Similar to solution 2, we need to maximise the equation $(x-1)(y-1)$. Starting from $x=y=223$, we simply do trial and error, adding $5$ to $x$ while subtracting $4$ from $y$ each time. We notice that the total number of basic squares increases as we increase $x$. 

 Thus, we simply need to find the first value for which the number of basic squares decrease or stay unchanged when we increase $x$, which occurs at $x=258$ and $y=195$, for which the total number of basic squares is $49,858$, which is less than when $x=253$ and $y=199$, which gives $49896$. Thus, our answer is $49896\mod1000 = \boxed{896}$.

 \href{https://mathproblemsolvingskills.wordpress.com/}{MathProblemSolvingSkills.com}

 


 




	\item (\textit{AIME II 2007}) \href{https://artofproblemsolving.com/wiki/index.php?title=Rectangle}{Rectangle} $ABCD$ is given with $AB=63$ and $BC=448.$ Points $E$ and $F$ lie on $AD$ and $BC$ respectively, such that $AE=CF=84.$ The \href{https://artofproblemsolving.com/wiki/index.php?title=Inscribed\_circle}{inscribed circle} of \href{https://artofproblemsolving.com/wiki/index.php?title=Triangle}{triangle} $BEF$ is \href{https://artofproblemsolving.com/wiki/index.php?title=Tangent}{tangent} to $EF$ at point $P,$ and the inscribed circle of triangle $DEF$ is tangent to $EF$ at \href{https://artofproblemsolving.com/wiki/index.php?title=Point}{point} $Q.$ Find $PQ.$


 
 \textbf{Answer: }259



\textbf{Solution}


                    \begin{figure}[H]
                        \centering
                        \includegraphics[width=0.5\textwidth]{Solution/AIME_II_2007_9_Solution_Figure_1.png}
                    \end{figure}
                    

 





\textbf{Solution 1}

Several \href{https://artofproblemsolving.com/wiki/index.php?title=Pythagorean\_triple}{Pythagorean triples} exist amongst the numbers given. $BE = DF = \sqrt{63^2 + 84^2} = 21\sqrt{3^2 + 4^2} = 105$. Also, the length of $EF = \sqrt{63^2 + (448 - 2\cdot84)^2} = 7\sqrt{9^2 + 40^2} = 287$.

 Use the \href{https://artofproblemsolving.com/wiki/index.php?title=Two\_Tangent\_Theorem}{Two Tangent Theorem} on $\triangle BEF$. Since both circles are inscribed in congruent triangles, they are congruent; therefore, $EP = FQ = \frac{287 - PQ}{2}$. By the Two Tangent theorem, note that $EP = EX = \frac{287 - PQ}{2}$, making $BX = 105 - EX = 105 - \left[\frac{287 - PQ}{2}\right]$. Also, $BX = BY$. $FY = 364 - BY = 364 - \left[105 - \left[\frac{287 - PQ}{2}\right]\right]$. 

 Finally, $FP = FY = 364 - \left[105 - \left[\frac{287 - PQ}{2}\right]\right] = \frac{805 - PQ}{2}$. Also, $FP = FQ + PQ = \frac{287 - PQ}{2} + PQ$. Equating, we see that $\frac{805 - PQ}{2} = \frac{287 + PQ}{2}$, so $PQ = \boxed{259}$.

 





\textbf{Solution 2}

By the \href{https://artofproblemsolving.com/wiki/index.php?title=Two\_Tangent\_Theorem}{Two Tangent Theorem}, we have that $FY = PQ + QF$. Solve for $PQ = FY - QF$. Also, $QF = EP = EX$, so $PQ = FY - EX$. Since $BX = BY$, this can become $PQ = FY - EX + (BY - BX)$$= \left(FY + BY\right) - \left(EX + BX\right) = FB - EB$. Substituting in their values, the answer is $364 - 105 = 259$.

 





\textbf{Solution 3}

Call the incenter of $\triangle BEF$ $O_1$ and the incenter of $\triangle DFE$ $O_2$. Draw triangles $\triangle O_1PQ,\triangle PQO_2$.

 Drawing $BE$, We find that $BE = \sqrt {63^2 + 84^2} = 105$. Applying the same thing for $F$, we find that $FD = 105$ as well. Draw a line through $E,F$ parallel to the sides of the rectangle, to intersect the opposite side at $E_1,F_1$ respectively. Drawing $\triangle EE_1F$ and $FF_1E$, we can find that $EF = \sqrt {63^2 + 280^2} = 287$. We then use Heron's formula to get:

 
 \[[BEF] = [DEF] = 11 466\].

 So the inradius of the triangle-type things is $\frac {637}{21}$.

 Now, we just have to find $O_1Q = O_2P$, which can be done with simple subtraction, and then we can use the \href{https://artofproblemsolving.com/wiki/index.php?title=Pythagorean\_Theorem}{Pythagorean Theorem} to find $PQ$. 

 





\textbf{Solution 4}

Why not first divide everything by its greatest common factor, $7$? Then we're left with much simpler numbers which saves a lot of time. In the end, we will multiply by $7$.

 From there, we draw the same diagram as above (with smaller numbers). We soon find that the longest side of both triangles is 52 (64 - 12). That means:

 $A = rs$ indicating $26(9)=r(54)$ so $r = 13/3$.

 Now, we can start applying the equivalent tangents. Calling them $a$, $b$, and $c$ (with $c$ being the longest and a being the shortest),

 $a+b+c$ is the semi perimeter or $54$. And since the longest side (which has $b+c$) is $52$, $a=2$.

 Note that the distance $PQ$ we desired to find is just $c - a$. What is $b$ then? $b = 13$. And $c$ is $39$. Therefore the answer is $37$... $NOT.$

 Multiply by $7$ back again (I hope you remembered to write this in $huge$ letters on top of the scrap paper!), we actually get $\boxed{259}$.

 





\textbf{Solution 5}

Scaling everything by 7, we have that $AE = 12, AB = 9, BF = 52$. Note that  if the perpendicular of $F$ dropped down to $ED$ is $X$, then $EX  = 52-12 = 40$. But $FX = 9$ and so we have a $9-40-41$ right triangle with $EFX$ meaning $EF = 41$. Now, by symmetry, we know that $EP = QF = a$ meaning $PF = 41-a$. If the tangent of the circle inscribed in $BEF$ is tangent to $BE$ at $Y$, then if $BY = b$ we have a system of equations. $a+b = 15, b+41-a = 52$. We can then solve for $a$, and since $PQ = 41-2a$, the rest follows.

 





\textbf{Solution 6 (Coordinate Geometry Bash)}

Set the origin at $D(0,0)$ with $A(0, 448)$, $B(63, 448)$, and $C(63, 0)$.Points $E$ and $F$ are defined by $AE=CF=84$:

 
                \begin{itemize}
                    \item $E = (0, 448 - 84) = (0, 364)$\item $F = (63, 84)$
                \end{itemize}


                The vertices of $\triangle DEF$ are $D(0,0)$, $E(0, 364)$, and $F(63, 84)$. The side lengths are:

 
                \begin{itemize}
                    \item $DE = 364$\item $DF = \sqrt{63^2 + 84^2} = 21\sqrt{3^2 + 4^2} = 105$\item $EF = \sqrt{(63-0)^2 + (84-364)^2} = \sqrt{63^2 + (-280)^2} = 7\sqrt{9^2 + 40^2} = 287$
                \end{itemize}


                Using the incenter formula $I = \left( \frac{ax_1 + bx_2 + cx_3}{a+b+c}, \frac{ay_1 + by_2 + cy_3}{a+b+c} \right)$, where $a, b, c$ are side lengths opposite to vertices $D, E, F$:

 
                \begin{itemize}
                    \item $a = EF = 287, \quad b = DF = 105, \quad c = DE = 364$\item Perimeter $P = 287 + 105 + 364 = 756$\item $I_1 = \left( \frac{287(0) + 105(0) + 364(63)}{756}, \frac{287(0) + 105(364) + 364(84)}{756} \right) = \left( \frac{22932}{756}, \frac{68796}{756} \right) = (30\frac{1}{3}, 91)$
                \end{itemize}


                Let $Q'$ be the tangency point of the incircle on side $DE$. Since $DE$ lies on the y-axis ($x=0$), $Q'$ is the projection of the incenter $I_1$ onto the y-axis, which is $(0, 91)$. By the properties of tangents from a vertex to an incircle, the distance from vertex $E$ to the tangency point $Q$ on segment $EF$ is equal to the distance from $E$ to $Q'$ on segment $DE$:

 
                \begin{itemize}
                    \item $EQ = EQ' = |364 - 91| = 273$
                \end{itemize}


                By the symmetry of the rectangle and the identical construction of the two triangles, the distance $FP$ in $\triangle BEF$ must equal the distance $EQ$ in $\triangle DEF$:

 
                \begin{itemize}
                    \item $FP = EQ = 273$
                \end{itemize}


                Since $P$ and $Q$ both lie on segment $EF$ (length $287$):
 \[PQ = |EF - (EQ + FP)| = |287 - (273 + 273)| = |287 - 546| = \boxed{259}\]

 




	\item (\textit{AIME II 2007}) Let $S$ be a \href{https://artofproblemsolving.com/wiki/index.php?title=Set}{set} with six \href{https://artofproblemsolving.com/wiki/index.php?title=Element}{elements}. Let $\mathcal{P}$ be the set of all \href{https://artofproblemsolving.com/wiki/index.php?title=Subset}{subsets} of $S.$ Subsets $A$ and $B$ of $S$, not necessarily distinct, are chosen independently and at random from $\mathcal{P}$. The \href{https://artofproblemsolving.com/wiki/index.php?title=Probability}{probability} that $B$ is contained in one of $A$ or $S-A$ is $\frac{m}{n^{r}},$ where $m$, $n$, and $r$ are \href{https://artofproblemsolving.com/wiki/index.php?title=Positive}{positive} \href{https://artofproblemsolving.com/wiki/index.php?title=Integer}{integers}, $n$ is \href{https://artofproblemsolving.com/wiki/index.php?title=Prime}{prime}, and $m$ and $n$ are \href{https://artofproblemsolving.com/wiki/index.php?title=Relatively\_prime}{relatively prime}. Find $m+n+r.$ (The set $S-A$ is the set of all elements of $S$ which are not in $A.$)


 
 \textbf{Answer: }710



\textbf{Solution 1}

Use \href{https://artofproblemsolving.com/wiki/index.php?title=Casework}{casework}:

 
                \begin{itemize}
                    <li>$B$ has 6 elements:
                \begin{itemize}
                    \item Probability: $\frac{1}{2^6} = \frac{1}{64}$\item $A$ must have either 0 or 6 elements, probability: $\frac{2}{2^6} = \frac{2}{64}$.
                \end{itemize}


                </li><li>$B$ has 5 elements:
                \begin{itemize}
                    \item Probability: ${6\choose5}/64 = \frac{6}{64}$\item $A$ must have either 0, 6, or 1, 5 elements. The total probability is $\frac{2}{64} + \frac{2}{64} = \frac{4}{64}$.
                \end{itemize}


                </li><li>$B$ has 4 elements:
                \begin{itemize}
                    \item Probability: ${6\choose4}/64 = \frac{15}{64}$\item $A$ must have either 0, 6; 1, 5; or 2,4 elements. If there are 1 or 5 elements, the set which contains 5 elements must have four emcompassing $B$ and a fifth element out of the remaining $2$ numbers. The total probability is $\frac{2}{64}\left({2\choose0} + {2\choose1} + {2\choose2}\right) = \frac{2}{64} + \frac{4}{64} + \frac{2}{64} = \frac{8}{64}$.
                \end{itemize}


                </li>
                \end{itemize}


                We could just continue our casework. In general, the probability of picking B with $n$ elements is $\frac{{6\choose n}}{64}$. Since the sum of the elements in the $k$th row of \href{https://artofproblemsolving.com/wiki/index.php?title=Pascal\%27s\_Triangle}{Pascal's Triangle} is $2^k$, the probability of obtaining $A$ or $S-A$ which encompasses $B$ is $\frac{2^{7-n}}{64}$. In addition, we must count for when $B$ is the empty set (probability: $\frac{1}{64}$), of which all sets of $A$ will work (probability: $1$).

 Thus, the solution we are looking for is $\left(\sum_{i=1}^6 \frac{{6\choose i}}{64} \cdot \frac{2^{7-i}}{64}\right) + \frac{1}{64} \cdot \frac{64}{64}$ $=\frac{(1)(64)+(6)(64)+(15)(32)+(20)(16)+(15)(8)+(6)(4)+(1)(2)}{(64)(64)}$ $=\frac{1394}{2^{12}}$$=\frac{697}{2^{11}}$.

 The answer is $697 + 2 + 11 = 710$.

 





\textbf{Solution 2}

we need $B$ to be a subset of $A$ or $S-A$we can divide each element of $S$ into 4 categories:

 
                \begin{itemize}
                    \item it is in $A$ and $B$\item it is in $A$ but not in $B$\item it is not in $A$ but is in $B$\item or it is not in $A$ and not in $B$
                \end{itemize}


                these can be denoted as $+A+B$, $+A-B$,$-A+B$, and $-A-B$

 we note that if all of the elements are in $+A+B$, $+A-B$ or $-A-B$we have that $B$ is a subset of $A$which can happen in $\dfrac{3^6}{4^6}$ ways

 similarly if the elements are in $+A-B$,$-A+B$, or $-A-B$we have  that $B$ is a subset of $S-A$which can happen in $\dfrac{3^6}{4^6}$ ways as well

 but we need to make sure we don't over-count ways that are in both sets these are when $+A-B$ or $-A-B$which can happen in $\dfrac{2^6}{4^6}$ waysso our probability is $\dfrac{2\cdot 3^6-2^6}{4^6}= \dfrac{3^6-2^5}{2^{11}}=\dfrac{697}{2^{11}}$.

 so the final answer is $697 + 2 + 11 = 710$.

 





\textbf{Solution 3}

$B$ must be in $A$ or $B$ must be in $S-A$. This is equivalent to saying that $B$ must be in $A$ or $B$ is disjoint from $A$. The probability of this is the sum of the probabilities of each event individually minus the probability of each event occurring simultaneously. There are $\binom{6}{x}$ ways to choose $A$, where $x$ is the number of elements in $A$. From those $x$ elements, there are ${2^x}$ ways to choose $B$. Thus, the probability that $B$ is in $A$ is the sum of all the values $\binom{6}{x}({2^x})$ for values of $x$ ranging from $0$ to $6$. For the second probability, the ways to choose $A$ stays the same but the ways to choose $B$ is now ${2^{6-x}}$. We see that these two summations are simply from the Binomial Theorem and that each of them is ${(2+1)^6}$. We subtract the case where both of them are true. This only happens when $B$ is the null set. $A$ can be any subset of $S$, so there are ${2^6}$ possibilities. Our final sum of possibilities is $2\cdot 3^6-2^6$. We have ${2^6}$ total possibilities for both $A$ and $B$, so there are ${2^{12}}$ total possibilities. This reduces down to $\dfrac{2\cdot 3^6-2^6}{4^6}= \dfrac{3^6-2^5}{2^{11}}=\dfrac{697}{2^{11}}$.The answer is thus $697 + 2 + 11 = 710$.

 





\textbf{Solution 4}

Let $|S|$ denote the number of elements in a general set $S$. We use complementary counting. 

 There is a total of $2^6$ elements in $P$, so the total number of ways to choose $A$ and $B$ is $(2^6)^2 = 2^{12}$.

 Note that the number of $x$-element subset of $S$ is $\binom{6}{x}$. In general, for $0 \le |A| \le 6$, in order for $B$ to be in neither $A$ nor $S-A$, $B$ must have at least one element from both $A$ and $S-A$. In other words, $B$ must contain any subset of $A$ and $S-A$ except for the empty set $\{\}$. This can be done in $\binom{6}{|A|} (2^{|A|} - 1)(2^{6-|A|} - 1)$ ways. As $|A|$ ranges from $0$ to $6$, we can calculate the total number of unsuccessful outcomes to be 
 \[\sum_{|A| = 0}^{6} \binom{6}{|A|} (2^{|A|} - 1)(2^{6-|A|} - 1) = 2702.\] So our desired answer is 
 \[1 - \dfrac{2702}{2^{12}} = \dfrac{697}{2^{11}} \Rightarrow \boxed{710}.\]

 -MP8148

 





\textbf{Solution 5}

To begin with, we note that there are $2^6$ subsets of $S$(which we can assume is $\{1,2,3,4,5,6\}$), including the null set. This gives a total of $(2^6)^2 = 2^{12}$ total possibilities for A and B. 

 Case 1: B is contained in A.If B has $0$ elements, which occurs in $\binom{6}{0}$ ways, A can be anything, giving us $\binom{6}{0} \cdot 2^6$. If B has $1$ element, A must contain that element, and then the remaining 5 are free to be in A or not in A. This gives us $\binom{6}{1} \cdot 2^5$. Summing, we end up with the binomial expansion of $(2 + 1)^6 = 3^6$.

 Case 2: B is contained in S-A.By symmetry, this case is the same as Case 1, once again giving us $3^6$ possibilities.

 Case 3: B is contained in both.We claim here that B can only be the null set. For contradiction, assume that there exists some element $x$ in B which satisfies this restriction. Then, A must contain $x$ as well, but we also know that $S-A$ contains $x$, contradiction. Hence, B is the null set, whereas A can be anything. This gives us $2^6$ possibilities.

 Since we have overcounted Case 3 in both of the other two cases, our final count is $2 \cdot 3^6 - 2^6$. This gives us the probability $\frac{2 \cdot 3^6 - 2^6}{2^{12}}$. Upon simplifying, we end up with $\frac{697}{2^{11}}$, giving the desired answer of $\boxed{710}$.- Spacesam

 


 \href{https://mathproblemsolvingskills.wordpress.com/}{MathProblemSolvingSkills.com}

 


 




	\item (\textit{AIME II 2007}) Two long \href{https://artofproblemsolving.com/wiki/index.php?title=Cylinder}{cylindrical} tubes of the same length but different \href{https://artofproblemsolving.com/wiki/index.php?title=Diameter}{diameters} lie \href{https://artofproblemsolving.com/wiki/index.php?title=Parallel}{parallel} to each other on a \href{https://artofproblemsolving.com/wiki/index.php?title=Plane}{flat surface}. The larger tube has \href{https://artofproblemsolving.com/wiki/index.php?title=Radius}{radius} $72$ and rolls along the surface toward the smaller tube, which has radius $24$. It rolls over the smaller tube and continues rolling along the flat surface until it comes to rest on the same point of its \href{https://artofproblemsolving.com/wiki/index.php?title=Circumference}{circumference} as it started, having made one complete revolution. If the smaller tube never moves, and the rolling occurs with no slipping, the larger tube ends up a \href{https://artofproblemsolving.com/wiki/index.php?title=Distance}{distance} $x$ from where it starts. The distance $x$ can be expressed in the form $a\pi+b\sqrt{c},$ where $a,$ $b,$ and $c$ are \href{https://artofproblemsolving.com/wiki/index.php?title=Integer}{integers} and $c$ is not divisible by the \href{https://artofproblemsolving.com/wiki/index.php?title=Square}{square} of any \href{https://artofproblemsolving.com/wiki/index.php?title=Prime}{prime}. Find $a+b+c.$


 
 \textbf{Answer: }179



\textbf{Solution}


                    \begin{figure}[H]
                        \centering
                        \includegraphics[width=0.5\textwidth]{Solution/AIME_II_2007_11_Solution_Figure_1.png}
                    \end{figure}
                    

 If it weren’t for the small tube, the larger tube would travel $144\pi$. Consider the distance from which the larger tube first contacts the smaller tube, until when it completely loses contact with the smaller tube.

 Drawing the radii as shown in the diagram, notice that the \href{https://artofproblemsolving.com/wiki/index.php?title=Hypotenuse}{hypotenuse} of the \href{https://artofproblemsolving.com/wiki/index.php?title=Right\_triangle}{right triangle} in the diagram has a length of $72 + 24 = 96$. The horizontal line divides the radius of the larger circle into $72 - 24 = 48$ on the top half, which indicates that the right triangle has leg of 48 and hypotenuse of 96, a $30-60-90 \triangle$. 

 Find the length of the purple arc in the diagram (the distance the tube rolled, but not the horizontal distance). The sixty degree \href{https://artofproblemsolving.com/wiki/index.php?title=Central\_angle}{central angle} indicates to take $\frac{60}{360} = \frac 16$ of the circumference of the larger circle (twice), while the $180 - 2(30) = 120^{\circ}$ central angle in the smaller circle indicates to take $\frac{120}{360} = \frac 13$ of the circumference. This adds up to $2 \cdot \frac 16 144\pi + \frac 13 48\pi = 64\pi$. 

 The actual horizontal distance it takes can be found by using the $30-60-90 \triangle$s. The missing leg is equal in length to $48\sqrt{3}$. Thus, the total horizontal distance covered is $96\sqrt{3}$. 

 Thus, we get $144\pi - 64\pi + 96\sqrt{3} = 80\pi + 96\sqrt{3}$, and our answer is $\boxed{179}$.

 \href{https://mathproblemsolvingskills.wordpress.com/}{MathProblemSolvingSkills.com}

 


 




	\item (\textit{AIME II 2007}) The increasing \href{https://artofproblemsolving.com/wiki/index.php?title=Geometric\_sequence}{geometric sequence} $x_{0},x_{1},x_{2},\ldots$ consists entirely of \href{https://artofproblemsolving.com/wiki/index.php?title=Integer}{integral} powers of $3.$ Given that


 $\sum_{n=0}^{7}\log_{3}(x_{n}) = 308$ and $56 \leq \log_{3}\left ( \sum_{n=0}^{7}x_{n}\right ) \leq 57,$


 find $\log_{3}(x_{14}).$


 
 \textbf{Answer: }091



\textbf{Solution 1}

Let $y_n = \log_3 (x_n)$. The second equation tells us that $y_7 = 56$, since the previous terms are negligible for the sum. 

 Then, since $\{y_n\}$ is an arithmetic sequence, we find 
 \[\sum_{i=0}^{7}y_i=4(y_0+56)=308 \implies y_0=21 \implies y_{14}=2y_7-y_0=\boxed{091}\]

 





\textbf{Solution 2}

Suppose that $x_0 = a$, and that the common \href{https://artofproblemsolving.com/wiki/index.php?title=Ratio}{ratio} between the terms is $r$. 

 
The first conditions tells us that $\log_3 a + \log_3 ar + \ldots + \log_3 ar^7 = 308$. Using the rules of \href{https://artofproblemsolving.com/wiki/index.php?title=Logarithm}{logarithms}, we can simplify that to $\log_3 a^8r^{1 + 2 + \ldots + 7} = 308$. Thus, $a^8r^{28} = 3^{308}$. Since all of the terms of the geometric sequence are integral powers of $3$, we know that both $a$ and $r$ must be powers of 3. Denote $3^x = a$ and $3^y = r$. We find that $8x + 28y = 308$. The possible positive integral pairs of $(x,y)$ are $(35,1),\ (28,3),\ (21,5),\ (14,7),\ (7,9),\ (0,11)$.

 
The second condition tells us that $56 \le \log_3 (a + ar + \ldots ar^7) \le 57$. Using the sum formula for a \href{https://artofproblemsolving.com/wiki/index.php?title=Geometric\_series}{geometric series} and substituting $x$ and $y$, this simplifies to $3^{56} \le 3^x \frac{3^{8y} - 1}{3^y-1} \le 3^{57}$. The fractional part $\approx \frac{3^{8y}}{3^y} = 3^{7y}$. Thus, we need $\approx 56 \le x + 7y \le 57$. Checking the pairs above, only $(21,5)$ is close.

 
Our solution is therefore $\log_3 (ar^{14}) = \log_3 3^x + \log_3 3^{14y} = x + 14y = \boxed{091}$.

 





\textbf{Solution 3}

All these integral powers of $3$ are all different, thus in base $3$ the sum of these powers would consist of $1$s and $0$s. Thus the largest value $x_7$ must be $3^{56}$ in order to preserve the givens. Then we find by the given that $x_7x_6x_5\dots x_0 = 3^{308}$, and we know that the exponents of $x_i$ are in an arithmetic sequence. Thus $56+(56-r)+(56-2r)+\dots +(56-7r) = 308$, and $r = 5$. Thus $\log_3 (x_{14}) = \boxed{091}$.

 





\textbf{Solution 4}

Like above, we use logarithmic identities to convert the problem into workable equations. We begin by labelling the powers of 3. Call $x_0$, $x_1$, $x_2$..., as $3^n$, $3^{n+m}$, and $3^{n+2m}$... respectively. With this format we can rewrite the first given equation as $n + n + m + n+2m + n+3m+...+n+7m = 308$. Simplify to get $2n+7m=77$. (1)

 Now, rewrite the second given equation as $3^{56} \leq \left( \sum_{n=0}^{7}x_{n} \right) \leq 3^{57}$. Obviously, $x_7$, aka $3^{n+7m}$ $<3^{57}$ because there are some small fractional change that is left over. This means $n+7m$ is $\leq56$. Thinking about the geometric sequence, it's clear each consecutive value of $x_i$ will be either a power of three times smaller or larger. In other words, the earliest values of $x_i$ will be negligible compared to the last values of $x_i$. Even in the best case scenario, where the common ratio is 3, the values left of $x_7$ are not enough to sum to a value greater than 2 times $x_7$ (amount needed to raise the power of 3 by 1). This confirms that $3^{n+7m} = 3^{56}$. (2)

 Use equations 1 and 2 to get $m=5$ and $n=21$. $\log_{3}{x_{14}} = \log_{3}{3^{21+14*5}} = 21+14*5 = \boxed{091}$

 
-jackshi2006

 





\textbf{Solution 5 (dum)}

Proceed as in Solution 3 for the first few steps. We have the sequence $3^{a},3^{a+n},3^{a+2n}...$. As stated above, we then get that $a+a+n+...+a+7n=308$, from which we simplify to $2a+7n=77$. From here, we just go brute force using the second statement (that $3^{56}\leq 3^{a}+...+3^{a+7n}\leq 3^{57}$). Rearranging the equation from earlier, we get
 \[n=11-\frac{2a}{7}\]from which it is clear that $a$ is a multiple of $7$. Testing the first few values of $a$, we get:Case 1 ($a=7,n=9$)The sequence is then $3^{7}+...+3^{70}$, which breaks the upper bound.Case 2 ($a=14,n=7$)The sequence is then $3^{14}+...+3^{63}$, which also breaks the upper bound.Case 3 ($a=21,n=5$)This is the first reasonable one, giving $3^{21}+...+3^{56}$. It seems like this would break the upper bound, but from some testing we get:
 \[3^{21}+...+3^{56}<3^{57}\]
 \[1+...+3^{35}<3^{36}\]
 \[1+...+3^{30}<2*3^{35}\]
 \[3^{5}+...+3^{30}<2*3^{35}-1\]
 \[1+...+3^{25}<2*3^{30}-3^{-5}\]Repeating this process over and over, we eventually get (while ignoring the extremely small fractions left over)
 \[1<2*3^{5}\]which confirms that this satisfies our upper bound. Thus $a=21,n=5$, so $x_14=3^{a+14n}\rightarrow3^{91}$. We then get the requested answer, $\log_3(3^{91})=\boxed{091}$ ~ amcrunner

 





\textbf{Solution 6}

Let $\log_3(x_n)=a_n$. Our goal is to find $a_{14}$. Since the sequence consists of integral powers of $3$, all $a_n$ are integers.

 From the first given condition, we have
 \[a_0+a_1+a_2+\cdots+a_7=308.\]

 From the second condition, we claim that $a_7=56$. Since all $x_n$ are powers of $3$, the largest possible values of the eight terms are $3^{56},3^{55},3^{54},\dots,3^{49}$. Even if these formed an infinite geometric series with first term $3^{56}$ and ratio $\frac13$, the sum would be $\frac{3^{57}}{2}$, whose base-$3$ logarithm lies strictly between $56$ and $57$. Since
 \[\sum_{n=0}^{7} x_n < \sum_{n=0}^{\infty} 3^{56}\left(\frac13\right)^n,\]it follows that $x_7=3^{56}$, so $a_7=56$.

 Let $d$ denote the common difference of the arithmetic sequence $\{a_n\}$. Then the first condition can be rewritten as
 \[56+(56-d)+(56-2d)+\cdots+(56-7d)=308.\]Simplifying,
 \[8\cdot 56 - 28d = 308,\]so
 \[-28d=-140 \quad \text{and} \quad d=5.\]

 Therefore, \(a_{14}=56+7\cdot 5=\boxed{091}.\)

 




	\item (\textit{AIME II 2007}) A \href{https://artofproblemsolving.com/wiki/index.php?title=Triangle}{triangular} \href{https://artofproblemsolving.com/wiki/index.php?title=Array}{array} of \href{https://artofproblemsolving.com/wiki/index.php?title=Square}{squares} has one square in the first row, two in the second, and in general, $k$ squares in the $k$th row for $1 \leq k \leq 11.$ With the exception of the bottom row, each square rests on two squares in the row immediately below (illustrated in given diagram). In each square of the eleventh row, a $0$ or a $1$ is placed. Numbers are then placed into the other squares, with the entry for each square being the sum of the entries in the two squares below it. For how many initial distributions of $0$'s and $1$'s in the bottom row is the number in the top square a \href{https://artofproblemsolving.com/wiki/index.php?title=Multiple}{multiple} of $3$?


 
                    \begin{figure}[H]
                        \centering
                        \includegraphics[width=0.5\textwidth]{Problem/AIME_II_2007_13_Problem_Figure_1.png}
                    \end{figure}
                    


 
 \textbf{Answer: }640



\textbf{Solution}

Label each of the bottom squares as $x_0, x_1 \ldots x_9, x_{10}$. 

 Through \href{https://artofproblemsolving.com/wiki/index.php?title=Induction}{induction}, we can find that the top square is equal to ${10\choose0}x_0 + {10\choose1}x_1 + {10\choose2}x_2 + \ldots {10\choose10}x_{10}$. (This also makes sense based on a combinatorial argument: the number of ways a number can "travel" to the top position going only up is equal to the number of times it will be counted in the final sum.)

 Examine the equation $\mod 3$. All of the coefficients from $x_2 \ldots x_8$ will be multiples of $3$ (since the numerator will have a $9$). Thus, the expression boils down to $x_0 + 10x_1 + 10x_9 + x_{10} \equiv 0 \mod 3$. Reduce to find that $x_0 + x_1 + x_9 + x_{10} \equiv 0 \mod 3$. Out of $x_0,\ x_1,\ x_9,\ x_{10}$, either all are equal to $0$, or three of them are equal to $1$. This gives ${4\choose0} + {4\choose3} = 1 + 4 = 5$ possible combinations of numbers that work. 

 The seven terms from $x_2 \ldots x_8$ can assume either $0$ or $1$, giving us $2^7$ possibilities. The answer is therefore $5 \cdot 2^7 = \boxed{640}$.

 \textbf{Note} The specific induction process: we want to prove that if the bottom row is $x_0, x_1 \cdots x_{n}$, the top square is ${n\choose0}x_0 + {n\choose1}x_1 + {n\choose2}x_2 + \ldots {n\choose{n}}x_{n}$.

 Base Case: when n = 1, this is obviously true.

 Induction Step: Assume that the statement is true for n, and we wish to prove it for $n + 1$. Therefore, we add another row below the $n + 1$ row, and name the squares as $a_0, a_1, \cdots a_{n+1}$. We obtain that $x_0 = a_0 + a_1, x_1 = a_1 + a_2, \cdots x_{n} = a_{n} + a_{n+1}$, and using the assumption, yield the top square equals ${n\choose0}(a_0 + a_1) + {n\choose1}(a_1 + a_2) + \ldots {n\choose{n}}(a_{n} + a_{n+1})$. While it's easy to prove that ${n\choose{k}} + {n\choose{k+1}} = {n+1\choose{k+1}}$ (just expand it), we find that the two equations are identical. 

 Note contributed by MVP\_Harry, friend me if you want ~

 Assume that the numbers filled in the 11th row is a b c d e f g h i j. Writing the number in each box, you will find that the coefficient of each letter turns out to be a pascal triangle.

 





\textbf{Solution 2}

If we label the bottom row of values $a_1,a_2,\ldots,a_11$ then the sum of the bottom row is $a_1+\ldots+a_11$, the sum of the second row is $a_1+2(a_2+\ldots+a_10)+a_11$ and this pattern continues until the top element is $(a_1+a_11)+2(a_2+a_10)+\ldots+4(a_5+a_7)+5a_6$. Then, the terms with coefficient 0 mod 3 can be anything, and we do casework on the other solutions mod 3. We get 640. ~joeythetoey

 




	\item (\textit{AIME II 2007}) Let $f(x)$ be a \href{https://artofproblemsolving.com/wiki/index.php?title=Polynomial}{polynomial} with real \href{https://artofproblemsolving.com/wiki/index.php?title=Coefficient}{coefficients} such that $f(0) = 1,$ $f(2)+f(3)=125,$ and for all $x$, $f(x)f(2x^{2})=f(2x^{3}+x).$ Find $f(5).$


 
 \textbf{Answer: }676



\textbf{Solution 1}

Let $r$ be a root of $f(x)$. Then we have $f(r)f(2r^2)=f(2r^3+r)$; since $r$ is a root, we have $f(r)=0$; therefore $2r^3+r$ is also a root. Thus, if $r$ is real and non-zero, $|2r^3+r|>r$, so $f(x)$ has infinitely many roots. Since $f(x)$ is a polynomial (thus of finite degree) and $f(0)$ is nonzero, $f(x)$ has no real roots.

 Note that $f(x)$ is not constant. We then find two complex roots: $r = \pm i$. We find that $f(i)f(-2) = f(-i)$, and that $f(-i)f(-2) = f(i)$. This means that $f(i)f(-i)f(-2)^2 = f(i)f(-i) \Longrightarrow f(i)f(-i)(f(-2)^2 - 1) = 0$. Thus, $\pm i$ are roots of the polynomial, and so $(x - i)(x + i) = x^2 + 1$ will be a factor of the polynomial. (Note: This requires the assumption that $f(-2)\neq1$. Clearly, $f(-2)\neq-1$, because that would imply the existence of a real root.)

 The polynomial is thus in the form of $f(x) = (x^2 + 1)g(x)$. Substituting into the given expression, we have 

 
 \[(x^2+1)g(x)(4x^4+1)g(2x^2)=((2x^3+x)^2+1)g(2x^3+x)\]
 \[(4x^6+4x^4+x^2+1)g(x)g(2x^2)=(4x^6+4x^4+x^2+1)g(2x^3+x)\]

 Thus either $4x^6+4x^4+x^2+1=(4x^4+1)(x^2+1)$ is 0 for any $x$, or $g(x)$ satisfies the same constraints as $f(x)$. Continuing, by infinite descent, $f(x) = (x^2 + 1)^n$ for some $n$.

 Since $f(2)+f(3)=125=5^n+10^n$ for some $n$, we have $n=2$; so $f(5) = \boxed{676}$.

 
Comment:The answer is clearly correct, but the proof has a gap, i.e. there is no reason that $f(-2)\neq1$.  Since $f(x)$ has no real roots, the degree must be even.  Consider $g(x)= f(x)/f(-x)$.  Then since $f$ is non-zero, $g(x)=g(2x^3+x)$. Now the function $2x^3+x$ applied repeatedly from some real starting value of x becomes arbitrarily large, and the limit of $g(x)$ as $|x|$ approaches infinity is 1, so $g(x)$=1 for all x, or $f(x)=f(-x)$.  Then $f(x)=h(x^2+1)$ for some polynomial $h(x)$, and $h(x^2+1)h(4x^4+1)=h(4x^6+4x^4+x^2+1) = h((x^2+1)(4x^4+1))$.  Now suppose h has degree m.  It is clearly monic.  Assume that the next highest non-zero coefficient in h is k.  Then, subtracting $((x^2+1)(4x^4+1))^m$ from both sides of the equation yields a polynomial equality with degree $4m+2k$ on the left and degree $6k$ on the right, a contradiction.  So $h(x)=x^m$, and $f(x)=(1+x^2)^m$.

 





\textbf{Solution * (Similar)}

Let $\alpha$ be a root of $f(x)$, so $f(\alpha) = 0$. Substituting $\alpha$ into the given functional equation $f(x)f(2x^2) = f(2x^3+x)$ yields:
 \[f(\alpha)f(2\alpha^2) = f(2\alpha^3+\alpha) \implies f(2\alpha^3+\alpha) = 0\]Thus, if $\alpha$ is a root, then $g(\alpha) = 2\alpha^3+\alpha$ must also be a root. To prevent the polynomial from having infinitely many distinct roots (which would force $f(x) \equiv 0$), the sequence of generated roots must remain finite and bounded.

 Taking magnitudes, the root generation map is $|\alpha_{\text{new}}| = |\alpha| \cdot |2\alpha^2 + 1|$. For the sequence of roots to not go out to infinity or go to $0$ (noting $f(0)=1 \neq 0$), we must look for trapped orbits. 

 Testing the boundary condition where both components satisfy $|\alpha| = 1$ and $|2\alpha^2 + 1| = 1$:
 \[|2\alpha^2+1|=1 \implies 2(\alpha\overline{\alpha})+\alpha^2+\overline{\alpha}^2=0 \implies \alpha^2+\overline{\alpha}^2=-2\] 
 \[\implies \Re (\alpha^2) = -1 \implies \alpha^2=-1 \implies \alpha = \pm i\]

 Since $g(i) = -i$ and $g(-i) = i$, these roots form a stable, closed finite 2-cycle orbit.

 Thus, $\pm i$ are the only valid roots, so we can write $f(x)$ in its factored form with multiplicity $n$:
 \[f(x) = a[(x-i)(x+i)]^n = a(x^2+1)^n\]Given $f(0) = 1$, we find $a(0^2+1)^n = 1 \implies a = 1$. Thus, $f(x) = (x^2+1)^n$.

 To solve for $n$, use the condition $f(2) + f(3) = 125$:
 \[(2^2+1)^n + (3^2+1)^n = 125\]
 \[5^n + 10^n = 125\]By inspection, $n=2$ is the unique solution since $5^2 + 10^2 = 25 + 100 = 125$.

 Finally, evaluate the polynomial at $x = 5$:
 \[f(5) = (5^2+1)^2 = 26^2 = \mathbf{676}\]

 





\textbf{Solution 2}

Let $r$ be a root of $f(x).$ This means that $f(r)f(2r^2)=f(2r^3+r).$ In other words, $2r^3+r$ is a root of $f(x)$ too. Since $f(x)$ can't have infinitely many roots, 
 \[Q(x)=P(P(\dotsb P(P(r)) \dotsb))\] is cyclic, where $P(x)=2x^3+x.$ Now, we will do casework.

 Case 1: $\deg f\geq1$

 Subcase 1: $|r|>1$

 This means that 
 \[|2r^3+r|\geq|2r^3|-|r|=|r|(2|r|^2-1)>|r|(2\cdot1^2-1)=|r|.\] It follows that $|2r^3+r|>|r|$ for all $r.$ This implies that $Q(r)$ can't be cyclic. Thus, it is impossible for $|r|>1$ to be true.

 Subcase 2: $|r|<1$

 This means that $|2r^3+r|\geq2|r^3|-|r|=|r|(|2r^2|-1)<|r|.$ It follows that $|2r^3+r|<|r|$ for all $r.$ This implies that $Q(r)$ can't be cyclic. Thus, it is impossible for $|r|>1$ to be true.

 Subcase 3: $|r|=1.$ 

 Since $|r|$ is not greater than or less than 1, $|r|=1.$ This means that all the roots of the polynomial have a magnitude of $1.$ More specifically, $|2r^3+r|$ has a magnitude of one. Since this would mean an equality condition from the triangle inequality, $2r^3$ and $r$ are collinear with the origin in the complex plane. In other words, $\frac{2r^3}{r}=\pm c\Leftrightarrow cr=2r^3\Leftrightarrow 2r^2=c\Leftrightarrow r=\pm\sqrt{\pm\frac{c}{2}},$ for some real constant $c.$ Now, from $|r|=1,$ we find that $\left|\pm\sqrt{\pm\frac{c}{2}}\right|=1\Leftrightarrow \sqrt{\pm\frac{c}{2}}=1\Leftrightarrow \pm\frac{c}{2}=1\Leftrightarrow c=\pm2.$ Putting this back into the equation, we find that $r=1,-1,i,-i.$ Now, this means that $2r^3+r=3,-3,i,-i.$ $3$ and $-3$ obviously doesn't have a magnitude of $1.$ Thus, $i,-i$ are the only possible roots of the polynomial. Since roots come in conjugate pairs, $f(x)=[(x-i)(x+i)]^n=(x^2+1)^n,$ works for all constants $n\neq0.$

 Case 2: $\deg f=0.$

 This means that $f(x)=c,$ for some constant $c.$ In other words, $c^2=c.$ We can easily find that this means that $c=0,1.$Combining all the cases, we conclude that $f(x)=(x^2+1)^n,0,1$ are the only polynomials that satisfy this equation.Now, we can test! $f(x)=0,1$ obviously don't satisfy $f(2)+f(3)=125.$ Thus, $f(x)=(x^2+1)^n.$ Substituting, we find that $5^n+10^n=125\Leftrightarrow n=2.$ We conclude that $f(5)=(5^2+1)^2=26^2=\boxed{676}.$

 





\textbf{Solution 3 (Hope and Cheese)}

We find $f(1)f(2)=f(3)$, so 
 \[(f(1)+1)f(2)=125\] In hopes of a clean polynomial representation of $f$, we assume that $f$ is an integer polynomial. We guess $f(1)=4, f(2)=25$, so $f(3)=100$. 

 Noting the squares, we let $g(x)=\sqrt{f(x)}$, so $g(0)=1, g(1)=2, g(2)=5, g(3)=10$. Checking finite differences, we find that $g$ is quadratic, and solving, we find $g(x)=x^2+1$. We can check that this does indeed satisfy the polynomial equation (thankfully).

 Thus, so does $f$, and our answer is 
 \[f(5)=g(5)^2=(5^2+1)^2=\boxed{676}\]

 \href{https://mathproblemsolvingskills.wordpress.com/}{MathProblemSolvingSkills.com}

 


 




	\item (\textit{AIME II 2007}) Four \href{https://artofproblemsolving.com/wiki/index.php?title=Circle}{circles} $\omega,$ $\omega_{A},$ $\omega_{B},$ and $\omega_{C}$ with the same \href{https://artofproblemsolving.com/wiki/index.php?title=Radius}{radius} are drawn in the interior of \href{https://artofproblemsolving.com/wiki/index.php?title=Triangle}{triangle} $ABC$ such that $\omega_{A}$ is \href{https://artofproblemsolving.com/wiki/index.php?title=Tangent}{tangent} to sides $AB$ and $AC$, $\omega_{B}$ to $BC$ and $BA$, $\omega_{C}$ to $CA$ and $CB$, and $\omega$ is \href{https://artofproblemsolving.com/wiki/index.php?title=Externally\_tangent}{externally tangent} to $\omega_{A},$ $\omega_{B},$ and $\omega_{C}$. If the sides of triangle $ABC$ are $13,$ $14,$ and $15,$ the radius of $\omega$ can be represented in the form $\frac{m}{n}$, where $m$ and $n$ are \href{https://artofproblemsolving.com/wiki/index.php?title=Relatively\_prime}{relatively prime} positive integers. Find $m+n.$


 
 \textbf{Answer: }389



\textbf{Solution 1 (Homothety)}


                    \begin{figure}[H]
                        \centering
                        \includegraphics[width=0.5\textwidth]{Solution/AIME_II_2007_15_Solution_Figure_1.png}
                    \end{figure}
                    

 First, apply \href{https://artofproblemsolving.com/wiki/index.php?title=Heron\%27s\_formula}{Heron's formula} to find that $[ABC] = \sqrt{21 \cdot 8 \cdot 7 \cdot 6} = 84$. The semiperimeter is $21$, so the \href{https://artofproblemsolving.com/wiki/index.php?title=Inradius}{inradius} is $\frac{A}{s} = \frac{84}{21} = 4$. 

 Now consider the \href{https://artofproblemsolving.com/wiki/index.php?title=Incenter}{incenter} $I$ of $\triangle ABC$. Let the \href{https://artofproblemsolving.com/wiki/index.php?title=Radius}{radius} of one of the small circles be $r$. Let the centers of the three little circles tangent to the sides of $\triangle{ABC}$ be $O_A$, $O_B$, and $O_C$. Let the center of the circle tangent to those three circles be $O$. Note that $\triangle O_A O_B O_C$ and $\triangle ABC$ are similar because they have parallel sides and have the same incenter. Also, since the inradius of $\triangle O_A O_B O_C$ is $4-r$ and the inradius of $\triangle ABC$ is $4,$ it follows that $\triangle O_A O_B O_C$ is the result of a \href{https://artofproblemsolving.com/wiki/index.php?title=Homothety}{homothety} $\mathcal{H}\left(I, \frac{4-r}{4}\right)$ on $\triangle ABC.$ Now, note that $OO_A = OO_B = OO_C = 2r$ are all circumradii of $\triangle O_A O_B O_C$, and the circumradius of $\triangle ABC$ is $R = \frac{abc}{4[ABC]} = \frac{65}{8}.$ From the ratio of the homothety, we have that $\frac{2r}{R}=\frac{4-r}{4}.$ Plugging in $R=\frac{65}{8}$ and solving yields $r = \frac{260}{129}$ and the final answer is $\boxed{389}$.

 
\href{https://latex.artofproblemsolving.com/9/4/7/947b7f06d947dbf8bc5d8f61cdd193c330377372.png}{https://latex.artofproblemsolving.com/9/4/7/947b7f06d947dbf8bc5d8f61cdd193c330377372.png}

 





\textbf{Solution 2}


                    \begin{figure}[H]
                        \centering
                        \includegraphics[width=0.5\textwidth]{Solution/AIME_II_2007_15_Solution_Figure_2.gif}
                    \end{figure}
                    

 Consider a 13-14-15 triangle.  $A=84.$  (By \href{https://artofproblemsolving.com/wiki/index.php?title=Heron\%27s\_Formula}{Heron's Formula} or by 5-12-13 and 9-12-15 right triangles.)

 The inradius is $r=\frac{A}{s}=\frac{84}{21}=4$, where $s$ is the semiperimeter.  Scale the triangle with the inradius by a linear scale factor, $u.$

 The circumradius is $R=\frac{abc}{4rs}=\frac{13\cdot 14\cdot 15}{4\cdot 4\cdot 21}=\frac{65}{8},$ where $a,$ $b,$ and $c$ are the side-lengths.  Scale the triangle with the circumradius by a \href{https://artofproblemsolving.com/wiki/index.php?title=Line}{linear} scale factor, $v$.

 Cut and combine the triangles, as shown.  Then solve for $4u$:

 
            \begin{enumerate}[label=(\roman*)]
                \item $\frac{65}{8}v=8u$
            \end{enumerate}
            
            \begin{enumerate}[label=(\roman*)]
                \item $v=\frac{64}{65}u$
            \end{enumerate}
            
            \begin{enumerate}[label=(\roman*)]
                \item $u+v=1$
            \end{enumerate}
            
            \begin{enumerate}[label=(\roman*)]
                \item $u+\frac{64}{65}u=1$
            \end{enumerate}
            
            \begin{enumerate}[label=(\roman*)]
                \item $\frac{129}{65}u=1$
            \end{enumerate}
            
            \begin{enumerate}[label=(\roman*)]
                \item $4u=\frac{260}{129}$
            \end{enumerate}
            The solution is $260+129=\boxed{389}$.

 





\textbf{Solution 3 (elementary)}

Let $A'$, $B'$, $C'$, and $O$ be the centers of circles $\omega_{A}$, $\omega_{B}$, $\omega_{C}$, $\omega$, respectively, and let $x$ be their radius.

 Now, triangles $ABC$ and $A'B'C'$ are similar by parallel sides, so we can find ratios of two quantities in each triangle and set them equal to solve for $x$.

 Since $OA'=OB'=OC'=2x$, $O$ is the circumcenter of triangle $A'B'C'$ and its circumradius is $2x$. Let $I$ denote the incenter of triangle $ABC$ and $r$ the inradius of $ABC$. Then the inradius of $A'B'C'=r-x$, so now we compute r. Computing the inradius by $A=rs$, we find that the inradius of $ABC$ is $4$. Additionally, using the circumradius formula $R=\frac{abc}{4K}$ where $K$ is the area of $ABC$ and $R$ is the circumradius, we find $R=\frac{65}{8}$. Now we can equate the ratio of circumradius to inradius in triangles $ABC$ and $A'B'C'$.

 
 \[\frac{\frac{65}{8}}{4}=\frac{2x}{4-x}\]

 Solving, we get $x=\frac{260}{129}$, so our answer is $260+129=\boxed{389}$.

 





\textbf{Solution 4}

According to the diagram, it is easily to see that there is a small triangle made by the center of three circles which aren't in the middle. The circumradius of them is$2r$. Now denoting $AB=13;BC=14;AC=15$, and centers of circles tangent to $AB,AC;AC,BC;AB,BC$ are relatively $M,N,O$ with $OJ,NK$ both perpendicular to $BC$. It is easy to know that $tanB=\frac{12}{5}$, so $tan\angle OBJ=\frac{2}{3}$ according to half angle formula. Similarly, we can find $tan\angle NCK=\frac{1}{2}$. So we can see that $JK=ON=14-\frac{7x}{2}$. Obviously, $\frac{2x}{14-\frac{7x}{2}}=\frac{65}{112}$ . After solving, we get$x=\frac{260}{129}$, so our answer is $260+129=\boxed{389}$. ~bluesoul

 \textbf{Sidenote (Generalization)} If four circles $\omega,$ $\omega_{A},$ $\omega_{B},$ and $\omega_{C}$ with the same radius are drawn in the interior of triangle $ABC$ such that $\omega_{A}$ is tangent to sides $AB$ and $AC$, $\omega_{B}$ to $BC$ and $BA$, $\omega_{C}$ to $CA$ and $CB$, and $\omega$ is externally tangent to $\omega_{A},$ $\omega_{B},$ and $\omega_{C}$. If $ABC$ has side lengths $a,b,$ and $c$, then the radius of $\omega$ can be written as 
 \[\frac{abc\sqrt{2b^{2}c^{2}+2a^{2}b^{2}+2a^{2}c^{2}-a^{4}-b^{4}-c^{4}}}{4b^{2}c^{2}+4a^{2}b^{2}+4a^{2}c^{2}-2a^{4}-2b^{4}-2c^{4}+2a^{2}bc+2ab^{2}c+2abc^{2}},\] or, more simply as, 
 \[\frac{abc\cdot K}{8K^{2}+sabc},\] where $K$ is the area of the triangle and $s$ is the semiperimeter

 





\textbf{Solution 5}


                    \begin{figure}[H]
                        \centering
                        \includegraphics[width=0.5\textwidth]{Solution/AIME_II_2007_15_Solution_Figure_3.png}
                    \end{figure}
                    

 Let $O_A, O_B, O_C, O$ be the centers of $w_A, w_B, w_C, w$, respectively. Also, let $I$ be the \href{https://artofproblemsolving.com/wiki/index.php?title=Incenter}{incenter} of $ABC$ and $r$ be the radius of circle $w$. Since $AB||O_AO_B$, $BC||O_BO_C$, and $CA||O_CO_A$, we know that 

 
 \[\angle BAI = \angle O_BO_AI, \angle CBI = \angle O_CO_BI, \angle ACI = \angle O_AO_CI \text{ and }\angle CAI = \angle O_CO_AI, \angle ABI = \angle O_AO_BI, \angle BCI = \angle O_BO_CI.\]

 That means $\angle ABC = \angle O_AO_BO_C$, $\angle BAC = \angle O_BO_AO_C$, and $\angle ACB = \angle O_AO_CO_B$. Thus, $\triangle ABC \sim \triangle O_AO_BO_C$. We also know that we are scaling each side of $\triangle ABC$ (from $AB$ to $O_AO_B$ for instance), about $I$ (since A,O\_A,I are \href{https://artofproblemsolving.com/wiki/index.php?title=Collinear}{collinear}; same apply with $B$ and $C$).

 Now, let the \href{https://artofproblemsolving.com/wiki/index.php?title=Homothety}{homothety} $\mathcal{H} (I, x)$ map $\triangle ABC$ to $\triangle O_AO_BO_C$. To start off, we know the \href{https://artofproblemsolving.com/wiki/index.php?title=Circumradius}{circumradius} of $O_AO_BO_C$ is $O$, since $OO_A = OO_B = OO_C = 2r$. Since $O_AO_B = 13x$, $O_BO_C = 14x$, $O_CO_A = 15x$, we can get an relationship involving $x$ and $r$ via another way to find the circumradius:

 
 \[[\triangle O_AO_BO_C ] =\frac{abc}{4R} \Longrightarrow 84x^2 =\frac{13x\cdot 14x\cdot 15x}{4\cdot 2r} \Longrightarrow r=\frac{65x}{16}\]

 Take notice of the \href{https://artofproblemsolving.com/wiki/index.php?title=Inradius}{inradius} of $ABC$. We get the inradius to be $[\triangle ABC ] = sr_0 \Longrightarrow r_0=4$. Let the tangency point of the \href{https://artofproblemsolving.com/wiki/index.php?title=Incircle}{incircle} and side $BC$ be $H$. We know $IH = 4$. We also know that we can cut off the part of $IH$ that is outside of $\triangle O_AO_BO_C$ to get the inradius of $\triangle O_AO_BO_C$. To part that is outside $\triangle O_AO_BO_C$ turns out just to be the radius of circle $w_B$ (as seen in the picture). That means the inradius of $\triangle O_AO_BO_C$ is just $4-r$. We can calculate that incradius in another way, though. We know that the inradius of $\triangle ABC$ is $4$, which means the inradius of $\triangle O_AO_BO_C$ is just $4x$ (by our homethety ratio).

 Thus, we have $4x = 4-r = 4-\dfrac{65x}{16} \Longrightarrow x = \dfrac{64}{129} \Longrightarrow r = \dfrac{260}{129}$. That gives $\boxed{389}$ as our final answer.

 The homethety turned out to be $\mathcal{H} \left(I, \dfrac{64}{129}\right)$

 





\textbf{Solution 6}

Let the radius of $\omega$ be $r$, let the centers of $\omega_A$, $\omega_B$, and $\omega_C$ be $O_A$, $O_B$, and $O_C$, respectively, and let the feet of the altitudes from $O_A$ and $O_C$ onto $AC$ be $X$ and $Y$, respectively. Now, observe that $\triangle O_AO_BO_C$ has all its sides parallel to the sides of $\triangle ABC$, so the two triangles are similar. Then, we get\begin{align*}AC&=AX+XY+YC,\\15&=r\cot\frac{\angle A}{2}+O_AO_C+r\cot\frac{\angle C}{2},\\15&=r\sqrt{\frac{1+\cos A}{1-\cos A}}+4r\sin\angle B+r\sqrt{\frac{1+\cos C}{1-\cos C}},\\15&=r\sqrt{\frac{65+33}{65-33}}+\frac{48}{13}r+r\sqrt{\frac{5+3}{5-3}},\\r&=\frac{15}{\frac 74+\frac{48}{13}+2}\\&=\boxed{\frac{260}{129}}.\,\square\\\end{align*}


 


 \href{https://mathproblemsolvingskills.wordpress.com/}{MathProblemSolvingSkills.com}

 


 




\end{enumerate}